\section{进程恢复、故障注入与负载测试}

\CaseHeader{FT-REC-001}{待建仓单期间重启调度器}{\OwnerAuto}{本地进程}{\StatusReady}

建立待建仓单后终止共享调度器，再由 Supervisor 拉起。预期根据稳定 client ID 恢复开放订单，不重复挂单；heartbeat、run ID 和 PID 更新。

\CaseHeader{FT-REC-002}{待平仓单期间重启调度器}{\OwnerAuto}{本地进程}{\StatusPass}

已有持仓和退出单时重启。预期从开放订单恢复 Cell；若数据库已有持仓但退出单尚未保存，启动后补建退出单。

\CaseHeader{FT-REC-003}{币安接受订单但响应丢失}{\OwnerAuto}{故障注入}{\StatusReady}

模拟 \texttt{POST /order} 已成功、客户端收到超时。若订单仍开放，应按 client ID 恢复；若订单已成交但无法查询，禁止重复建仓并将真实仓位标记为未分配。该用例是迁移前最高优先级之一。

\CaseHeader{FT-REC-004}{SQLite 写入失败}{\OwnerAuto}{故障注入}{\StatusReady}

在订单创建前后分别注入数据库异常，验证不会悄悄丢失错误。恢复后应通过开放订单/client ID 或资源池进入可审计状态，不得无上限重试下单。

\CaseHeader{FT-REC-005}{持仓/订单接口超时与 429/5xx}{\OwnerAuto}{故障注入}{\StatusReady}

模拟超时、限流和服务端错误。预期本轮跳过危险重写，错误写入 heartbeat；后续成功轮次恢复。不得对 429 进行无间隔紧密重试。

\CaseHeader{FT-REC-006}{整机重启}{\OwnerHuman+\OwnerLive}{服务器操作}{\StatusManual}

人工在测试服务器重启服务或整机，Codex 检查单调度器 PID、SQLite WAL、开放订单恢复和未重复下单。此用例必须在小额测试策略上执行。

\CaseHeader{FT-LOAD-001}{50 组乘 5 Cell 单币对}{\OwnerAuto}{模拟负载}{\StatusPass}

验证 250 个 Cell 共用一次 mark price、一次 symbol filters 和一次开放订单快照；client ID 全部唯一，未到轮询时间不访问交易所。

\CaseHeader{FT-LOAD-002}{50 币对乘 5 Cell}{\OwnerAuto}{模拟负载}{\StatusPass}

验证请求量按币对数而不是 Cell 数增长；同轮每币对最多一次价格和开放订单快照。记录执行时间和 Python 峰值内存。

\CaseHeader{FT-LOAD-003}{2U2GB 服务器资源观察}{\OwnerHuman+\OwnerLive}{服务器资源}{\StatusManual}

部署 50 组低频测试配置，观察 24 小时 CPU、RSS、SQLite 文件增长、API 错误和 Binance 请求权重。验收建议：无 OOM、无重复调度器、无持续 429、单轮协调完成时间明显小于最短协调间隔。

\CaseHeader{FT-SOAK-001}{24 小时小额稳定性}{\OwnerHuman+\OwnerLive}{真实资金}{\StatusManual}

先运行 3--5 组，每组 1--3 Cell，轮询间隔 1--3 分钟。人工每日检查一次币安仓位与 Web 资源池，Codex 汇总事件和异常。24 小时无严重问题后才进入 72 小时与更大组数。
